\chapter{Отношения и пропорции}\label{chap:c05}\markright{Глава 5}
\begin{definition}\term{Отношение} (\emph{ratio}) \(a\) к \(b\) — частное \(\frac{a}{b}\) при \(b \ne  0\); обозначается \(a : b\).\end{definition}
\begin{definition}\term{Пропорция} (\emph{proportion}) — равенство двух отношений: \(\frac{a}{b} = \frac{c}{d}\), где \(b,\allowbreak d \ne  0\).\end{definition}
\begin{definition}\term{Прямая пропорциональность} (\emph{direct proportion}) — зависимость \(y = kx\) с постоянным ненулевым коэффициентом \(k\).\end{definition}
\begin{definition}\term{Обратная пропорциональность} (\emph{inverse proportion}) — зависимость \(y = \frac{k}{x}\) при \(x \ne  0\) с постоянным ненулевым \(k\).\end{definition}
\begin{definition}\term{Масштаб} (\emph{scale}) \(1 : N\) — отношение длины на карте к длине на местности в одинаковых единицах.\end{definition}
\section{Основное свойство пропорции}
При \(b,\allowbreak d \ne  0\) равенство \(\frac{a}{b} = \frac{c}{d}\) равносильно \(a \cdot  d = b \cdot  c\).\par
\section{Правила применения}
При прямой пропорциональности увеличение \(x\) в несколько раз увеличивает \(y\) во столько же раз.\par
При обратной пропорциональности положительных величин увеличение \(x\) в несколько раз уменьшает \(y\) во столько же раз.\par
Деление \(S\) в отношении \(m : n\) при \(m,\allowbreak n > 0\): части равны \(S \cdot  \frac{m}{m + n}\) и \(S \cdot  \frac{n}{m + n}\).\par
{\small\textit{Пример: \(60\) разделить в отношении \(2 : 3\) — это \(24\) и \(36\), поскольку одна часть равна \(\frac{60}{5} = 12\).}}\par
\Needspace{7\baselineskip}\section*{Задачи}
\begingroup\setlength{\parskip}{1pt}
\noindent\begin{minipage}[t]{\linewidth}\textbf{5.1.} Сократи отношение \(18 : 24\).\par\end{minipage}\par
\noindent\begin{minipage}[t]{\linewidth}\textbf{5.2.} Реши пропорцию \(\dfrac{x}{15} = \dfrac{4}{5}\).\par\end{minipage}\par
\noindent\begin{minipage}[t]{\linewidth}\textbf{5.3.} Реши пропорцию \(\dfrac{7}{x} = \dfrac{21}{18}\).\par\end{minipage}\par
\noindent\begin{minipage}[t]{\linewidth}\textbf{5.4.} Раздели \(64\) в отношении \(3 : 5\).\par\end{minipage}\par
\noindent\begin{minipage}[t]{\linewidth}\textbf{5.5.} Для \(4\) порций нужно \(300\) г крупы. Сколько нужно для \(10\) порций?\par\end{minipage}\par
\par\addvspace{5pt}\begin{lrbox}{\CatFlowBox}\begin{minipage}[t]{\linewidth}
\setlength{\parskip}{2pt}\setlength{\columnsep}{4.5mm}\setlength{\intextsep}{0pt}
\begin{wrapfigure}[12]{r}{54mm}
\includegraphics[width=54mm,height=58mm,keepaspectratio]{images/unique/art-21.png}
\label{catflow:21}
\end{wrapfigure}
\textbf{5.6.} \(5\) одинаковых тетрадей стоят \(180\) рублей. Сколько стоят \(8\)?\par

\textbf{5.7.} На карте масштаба \(1 : 100 000\) расстояние \(3,5\) см. Найди расстояние на местности в километрах.\par

\textbf{5.8.} \(6\) рабочих выполняют работу за \(10\) дней. За сколько дней справятся \(15\) рабочих при той же производительности?\par

\textbf{5.9.} Поездка при скорости \(60\) км/ч занимает \(3\) часа. Сколько займёт при \(90\) км/ч на том же пути?\par

\textbf{5.10.} Смешивают сок и воду в отношении \(2 : 3\). Сколько каждого компонента нужно для \(1{,}5\) л смеси?\par

\textbf{5.11.} При \(y = 3 x\) найди \(y\) для \(x = 4\). Как изменится \(y\) при удвоении \(x\)?\par

\textbf{5.12.} При \(y = \dfrac{24}{x}\) найди \(y\) для \(x = 6\). Как изменится \(y\) при удвоении положительного \(x\)?\par
\par\WFclear\end{minipage}\end{lrbox}%
\ifdim\dimexpr\ht\CatFlowBox+\dp\CatFlowBox\relax<64mm\dp\CatFlowBox=\dimexpr64mm-\ht\CatFlowBox\relax\fi%
\noindent\usebox{\CatFlowBox}\par\addvspace{5pt}

\noindent\begin{minipage}[t]{\linewidth}\textbf{5.13.} Для \(6\) порций супа нужно \(450\) г картофеля. Сколько нужно для \(10\) таких же порций?\par\end{minipage}\par
\noindent\begin{minipage}[t]{\linewidth}\textbf{5.14.} \(4\) одинаковых насоса заполняют бассейн за \(9\) часов. За сколько часов его заполнят \(6\) таких насосов, работающих одновременно?\par\end{minipage}\par
\noindent\begin{minipage}[t]{\linewidth}\textbf{5.15.} Для напитка сироп смешивают с водой в отношении \(1 : 7\). Сколько сиропа и воды нужно для \(2\) л напитка?\par\end{minipage}\par
\noindent\begin{minipage}[t]{\linewidth}\textbf{5.16.}\starmark{} Рабочие одинаковой производительности: \(6\) человек выполняют работу за \(8\) дней по \(5\) часов. За сколько дней справятся \(4\) человека по \(6\) часов?\par\end{minipage}\par
\noindent\begin{minipage}[t]{\linewidth}\textbf{5.17.}\starmark{} В смеси отношение сиропа к воде \(2 : 5\). После добавления \(3\) л воды стало \(1 : 4\). Сколько литров смеси было сначала?\par\end{minipage}\par
\noindent\begin{minipage}[t]{\linewidth}\textbf{5.18.}\starmark{} Автомобиль ехал одинаковое время со скоростями \(60\) и \(90\) км/ч. Найди среднюю скорость. А если одинаковыми были расстояния?\par\end{minipage}\par
\noindent\begin{minipage}[t]{\linewidth}\textbf{5.19.} Пусть \(b,d>0\) и \(\dfrac ab=\dfrac cd=k\). Докажи, что \(\dfrac{a+c}{b+d}=k\).\par\end{minipage}\par
\noindent\begin{minipage}[t]{\linewidth}\textbf{5.20.}\starmark{} Пусть \(a,b,c,d>0\) и \(\dfrac ab<\dfrac cd\). Докажи, что \(\dfrac ab<\dfrac{a+c}{b+d}<\dfrac cd\).\par\end{minipage}\par
\endgroup
